\section{Collar and radius scales}

\begin{definition}[Radius exponent]
    \label{def:area_law_radius_exponent}
    \lean{TNLean.PEPS.AreaLaw.Exponents.radiusExponent}
    \uses{thm:area_law_exponent_gaps}
    For the fixed amplification parameters, let
    $\gamma=(1-\varepsilon)/\alpha$.
    This is the exponent in the radius bound of
    \cite[Proposition~10.2]{OpenAI2026AreaLaw}.
\end{definition}

\begin{theorem}[Separation of the collar scales]
    \label{thm:area_law_radius_exponent_gaps}
    \lean{TNLean.PEPS.AreaLaw.Exponents.radius_exponent_gaps,
        TNLean.PEPS.AreaLaw.Exponents.collar_radius_scale_tendsto_zero}
    \uses{def:area_law_radius_exponent}
    The fixed parameters satisfy $0<\gamma$ and
    $1-\varepsilon<\gamma<\beta$. For every real constant $C$,
    $[L^{1-\varepsilon}+2CL^\gamma]/L^\beta\to0$
    as the positive integer $L$ tends to infinity.
\end{theorem}
\begin{proof}
    The inequalities follow by substituting the fixed parameters.
    Each power in the numerator has exponent strictly smaller than $\beta$.
\end{proof}

\begin{theorem}[Uniform control of rounded collars]
    \label{thm:area_law_uniform_collar_threshold}
    \lean{TNLean.PEPS.AreaLaw.Exponents.exists_collar_radius_threshold,
        TNLean.PEPS.AreaLaw.Exponents.collar_add_radius_isLittleO}
    \uses{thm:area_law_radius_exponent_gaps}
    For every $C\in\mathbb R$ and $\eta>0$, there is $N\ge2$ such that
    every integer $L\ge N$ and every nonnegative integer radius
    $r\le CL^\gamma$ satisfy
    $\lfloor L^{1-\varepsilon}\rfloor+2r<\eta L^\beta$.
    Consequently, for every sequence of nonnegative integer radii satisfying
    $r(L)\le CL^\gamma$ eventually,
    $\lfloor L^{1-\varepsilon}\rfloor+2r(L)=o(L^\beta)$.
    The threshold depends only on $C$ and $\eta$.
\end{theorem}
\begin{proof}
    Apply the preceding limit and use
    $\lfloor L^{1-\varepsilon}\rfloor\le L^{1-\varepsilon}$.
    This proves the scale comparison in the final paragraph of the proof of
    \cite[Proposition~10.2]{OpenAI2026AreaLaw}; the geometric clearance and
    amplified mutual-information estimate require separate arguments.
\end{proof}
